% Sources:
% - docs/paper/literature/citation-verification.tsv
% - docs/paper/literature/literature-corpus.md
% - docs/paper/literature/novelty-positioning.md
% - docs/paper/literature/related-work-revision-plan.md
% - docs/paper/related-work-draft.md
% - docs/paper/related-work-revision-report.md

\subsection{K-Waay and Post-Quantum Asynchronous Key Exchange}

\kwaay is the direct protocol source for this work. Collins et al. construct a
post-quantum, X3DH-like deniable authenticated key exchange using split-KEM
techniques, introduce \batchreceive for receiver-side processing with reused
ephemeral material, and provide computational security definitions and proofs
for their construction \cite{collins2024kwaay}. The same specification states
that the elements of one \batchreceive call correspond to different parties
\cite[Sec.~5.1, p.~21]{collins2024kwaayfull}. Thus, the distinct-party condition is prior
work: this paper neither discovers it nor proposes a repair for an omitted
check. Our question is instead what party-level relation the stated condition
preserves when isolated at the batch-admission boundary.

Formal analyses of asynchronous secure messaging establish a broader but
different context. The analysis of Signal models X3DH and the Double Ratchet
as a multi-stage authenticated key-exchange system and proves scoped
cryptographic-session guarantees \cite{cohngordon2020signal}. Bhargavan et al.
translate the PQXDH specification into ProVerif and CryptoVerif models,
clarifying specification assumptions while analyzing authentication,
confidentiality, post-quantum forward secrecy, and co-deployment
\cite{bhargavan2024pqxdh}. The X3DH specification separately discusses
authentication, replay, key reuse, and identity binding
\cite{marlinspike2016x3dh}. These works show the importance of precise
identity and session modeling, but their principal objects are two-party
handshakes, ratcheting sessions, and cryptographic key security. Our model does
not reverify those constructions; it isolates a same-batch party relation and
ends at a symbolic receiver-acceptance event.

\subsection{Formal Analysis of Secure-Messaging Composition}

Prior work has already established that composition layers can carry security
obligations not implied by the guarantees of their components. Cremers,
Jacomme, and Naska analyze an application that manages and merges multiple
ratcheting sessions into a conversation, showing why conversation-level
post-compromise security must be studied separately from the security of an
individual chain \cite{cremers2023session}. This is the closest methodological
precedent for treating an interface above a secure-messaging primitive as an
independent formal object. The composition object is nevertheless different:
their work concerns session-to-conversation management, whereas ours concerns
party projections among entries of one receiver batch.

Group-messaging research gives further examples of identity-bearing
composition without defining the invariant studied here. TreeSync provides an
executable formal account of authenticated MLS group-management state and its
consistency and integrity properties \cite{wallez2023treesync}. Cryptographic
administration formalizes authorization and membership consistency for group
updates \cite{balbas2023administration}, while modular MLS analysis derives
secure-group-messaging guarantees from explicitly separated cryptographic
components \cite{alwen2021modular}. Cross-group work additionally shows that
healing properties depend on how users and compromise effects are represented
across groups \cite{cremers2021healing}. These results concern evolving group
state, membership authorization, cryptographic component composition, or
cross-group compromise. None is equivalent to requiring two positions in one
\batchreceive batch to project to different modeled parties.

The broader secure-messaging taxonomy likewise distinguishes participant
consistency from speaker, destination, authentication, and message-oriented
properties \cite{unger2015sok}. Participant consistency concerns agreement
among honest parties on their view of the participant list; our predicate
instead concerns party projections at different positions inside one batch.
These are different relations. General compositional-verification frameworks
formalize conditions under which component properties are preserved when
protocols are combined \cite{andova2008compositional}. Together, this
literature makes composition-level reasoning established prior art. Our use of
Tamarin's multiset-rewriting trace framework is also inherited methodology
\cite{meier2013tamarin}; the tool does not by itself enlarge the fixed
two-slot abstraction. The present contribution is limited to a different
composition object and property: same-batch party distinctness at a specific
\kwaay interface.

\subsection{Identity Semantics and Misbinding}

The security relevance of principal names and identity interpretation is also
well established. Lupetti, Dillema, and Stabell-Kul\o{} argue that names must
remain unique in an appropriate correlated-session scope to retain their
intended protocol meaning, and discuss local and global ways to verify that
assumption \cite{lupetti2006names}. Chosen-name analysis shows, from another
direction, that an adversary's ability to select or assign agent names changes
the formal attack surface \cite{ceelen2008chosenname}. These works preclude a
broad novelty claim for identity uniqueness or for treating identity
representation as a modeling decision.

Our property uses a narrower relation. It assumes an explicit party coordinate
and asks whether the same coordinate may occupy two slots of one admitted
batch. It does not analyze whether two principals share an ambiguous name, nor
does it study adversarial synthesis of party identifiers. Applying the
party/message/session/slot distinction to \batchreceive is therefore a
protocol-specific use of established identity semantics, not a new general
theory of identity.

Misbinding and unknown key-share work is conceptually adjacent but defines a
different failure. Formal misbinding analyses study cases in which an endpoint,
device, user, or protocol run becomes associated with the wrong identity
\cite{peltonen2020misbinding}. The TLS/SDP unknown key-share analysis similarly
concerns binding an external identity to another endpoint's key, fingerprint,
or session \cite{thomson2021uks}. Our repeated-party composition instead uses
the same modeled party coordinate twice; the current model proves neither an
incorrect endpoint-to-identity association nor peer disagreement about a
shared key. We therefore do not classify its witness as a misbinding, device
pairing attack, or unknown key-share attack.

\subsection{Injectivity, Replay, and Coordinate-Level Controls}

Injective agreement is an established authentication notion. Lowe's hierarchy
distinguishes aliveness, weak agreement, non-injective agreement, and
injective agreement, with injectivity adding a one-to-one relation between the
relevant protocol runs \cite{lowe1997hierarchy}. Sattarzadeh and Fallah give an
automated type-based treatment of injective agreement and use unique matching
begin/end resources to exclude duplicated receiver occurrences supported by
one transmitted message \cite{sattarzadeh2015injective}. Replay and repeated
acceptance are likewise explicit protocol concerns in the X3DH specification
\cite{marlinspike2016x3dh}. Neither injectivity nor replay-sensitive
correspondence is introduced by this paper.

We use scoped occurrence injectivity only as a diagnostic control. As reported
in Section~5, within one batch and receiver context the message-level variant
permits at most one acceptance occurrence for each exact sender-origin tuple,
while still admitting two different messages from the same party. This does
not guarantee delivery of every Send or exclude reuse across batches. It does
not show that existing
injectivity notions are inadequate; it shows that occurrence correspondence
and party distinctness state different relations. Similarly, exact-message
deduplication addresses equality of the message coordinate. The controlled
model supports only the narrower conclusion that such a restriction does not
imply the party-level objective under the encoded lifecycle. It does not imply
that replay protection or deduplication is ineffective for its own purpose.

\subsection{Positioning of This Work}

The closest works discussed above address the protocol condition itself
\cite{collins2024kwaayfull}, the composition-layer precedent
\cite{cremers2023session}, the established semantics of scoped name uniqueness
\cite{lupetti2006names}, and the occurrence-level injectivity vocabulary
\cite{sattarzadeh2015injective}. Their analysis objects differ from the
specific K-Waay batch-admission comparison studied here. This comparison,
rather than a global priority claim for identity or composition reasoning,
defines the scope of our contribution.

Against that background, two aspects of the present analysis have direct
support within the stated boundary. First, its auxiliary message-level
control exhibits a separation in which message distinction and scoped
occurrence injectivity hold while a same-party/different-message composition
remains reachable. Second, it places this separation between a relaxed
receiver-event witness and a non-vacuous party-level realization of the stated
\batchreceive condition. Both claims concern the fixed two-slot admission abstraction and
its recorded properties, not arbitrary batch lengths or complete \kwaay
executions.

Other aspects require correspondingly modest positioning. The distinction
between messages and identities, the general significance of name binding,
injective correspondence, and composition-level reasoning all predate this
work. Likewise, separating a semantic invariant from one possible enforcement
mechanism is not presented as a new general framework. The contribution is to
apply those established distinctions to the specific \batchreceive
party-admission question and to compare a plausible coordinate-level control
with direct preservation of the stated party relation.

This positioning makes no claim that the original protocol omits its stated
condition, that its computational theorem is invalid, or that a deployed
implementation violates the invariant. It also does not infer a cryptographic
break, misbinding, unknown key-share attack, application-level consequence, or
an arbitrary-size theorem from the bounded receiver-event evidence.
